\chapter{Algemene gravitatie en gewicht}\label{ch:g10:universal-gravitation}

Laat een appel los: hij valt. De Maan, zestig aardstralen hoger, lijkt
dat nooit te doen. Newtons grote inzicht is dat de Maan wel degelijk
\emph{valt} --- ongeveer 1.4 millimeter per seconde, precies zoveel als
een aantrekking die afneemt met het kwadraat van de afstand voorspelt.
Eén formule, met één universele constante, regeert de appel, de Maan, de
getijden en de planeten. Het onderbouwvolume beschreef de \omterm{def:g10:universal-gravitation:force}{gravitatie} met
woorden; dit hoofdstuk meet haar.

\section{De wet van de algemene gravitatie}

\begin{definition}[Gravitatiekracht]\label{def:g10:universal-gravitation:force}
Elke twee lichamen trekken elkaar aan, waaruit ze ook bestaan en hoe ver
ze ook uit elkaar liggen. Die aantrekking is de
\emph{gravitatiekracht}\index{gravitatiekracht}, en het verschijnsel heet
\emph{gravitatie}\index{gravitatie}. Ze heeft geen contact, geen touw en
geen tussenstof nodig, en ze is \emph{wederzijds}: elk lichaam trekt aan
het andere.
\end{definition}

\begin{theorem}[Wet van de algemene gravitatie]\label{thm:g10:universal-gravitation:law}
Twee lichamen met massa's $m_1$ en $m_2$ (in \unit{kg}), waarvan de
middelpunten een afstand $d$ uit elkaar liggen (in \unit{m}), trekken
elkaar aan met twee krachten langs de verbindingslijn van de
middelpunten, elk gericht naar het andere lichaam, met dezelfde grootte
\[
F = G\,\frac{m_1 m_2}{d^2}
\qquad \text{(in \unit{N})},
\]
waarin $G$ voor elk paar lichamen in het heelal dezelfde constante is.
\end{theorem}
\admitted

\begin{remark}[Waar de wet vandaan komt]
Newton destilleerde deze wet uit de waargenomen beweging van de planeten
(de wetten van Kepler); de eerlijke afleiding staat in het volume van
bachelorjaar 1. \Cref{ch:g12:satellites-kepler} zet de wet al aan het
werk op \omterm{def:g10:universal-gravitation:satellite}{satellieten} en planeten.
\end{remark}

\begin{definition}[De gravitatieconstante]\label{def:g10:universal-gravitation:constant}
De constante
\[
G = \qty{6.67e-11}{N.m^2/kg^2}
\]
heet de \emph{gravitatieconstante}\index{gravitatieconstante}. Haar
minuscule waarde verklaart waarom de zwaartekracht tussen alledaagse
voorwerpen onmerkbaar is: alleen wanneer minstens één massa astronomisch
is, stelt $G m_1 m_2 / d^2$ iets voor.
\end{definition}

\begin{omfigure}
\begin{tikzpicture}[scale=1.0]
  \draw[omExo, dashed] (0,0) -- (5.2,0);
  \fill[omDef!25] (0,0) circle (0.45);
  \draw[omDef, thick] (0,0) circle (0.45);
  \fill[omDef!25] (5.2,0) circle (0.28);
  \draw[omDef, thick] (5.2,0) circle (0.28);
  \node[font=\small] at (0,0) {$m_1$};
  \node[font=\small] at (5.2,0) {$m_2$};
  \draw[-Stealth, omThm, very thick] (0.45,0) -- (1.55,0)
    node[midway, above, font=\small] {$\vect F_{2/1}$};
  \draw[-Stealth, omThm, very thick] (4.92,0) -- (3.82,0)
    node[midway, above, font=\small] {$\vect F_{1/2}$};
  \draw[Stealth-Stealth, thin] (0,-0.75) -- (5.2,-0.75)
    node[midway, below, font=\small] {$d$};
\end{tikzpicture}

{\small De twee krachten van een gravitatiepaar: langs de
verbindingslijn van de middelpunten, naar het andere lichaam toe, met
\emph{dezelfde} grootte $F = G m_1 m_2 / d^2$, ook als de massa's
volstrekt ongelijk zijn.}
\end{omfigure}

\begin{remark}[De formule lezen]\label{rem:g10:universal-gravitation:reading}
Vier eigenschappen verdienen aandacht.
\begin{itemize}
  \item \emph{Wederzijds en even groot.} De Aarde trekt aan jou met jouw
        \omterm{def:g10:universal-gravitation:weight}{gewicht} --- en jij trekt aan de Aarde met precies dezelfde
        kracht. (De Aarde reageert er alleen minder op: ze is ongeveer
        $\num{e23}$ keer zwaarder.)
  \item \emph{Omgekeerd kwadratisch.} $d$ verdubbelen deelt de kracht
        door $4$; tien keer verder gaan deelt haar door $100$.
  \item \emph{Afstand tussen de middelpunten.} Voor bolvormige lichamen
        --- sterren, planeten, kanonskogels --- meet je $d$ van
        middelpunt tot middelpunt, niet van oppervlak tot oppervlak.
  \item \emph{Beide massa's tellen.} Een van beide massa's verdubbelen
        verdubbelt de kracht; de formule is symmetrisch in $m_1$ en
        $m_2$.
\end{itemize}
\end{remark}

\begin{example}[Twee vrienden]\label{ex:g10:universal-gravitation:friends}
Twee vrienden van \qty{70}{kg} staan met hun middelpunten \qty{1.0}{m}
uit elkaar:
\[
F = \num{6.67e-11} \times \frac{70 \times 70}{1.0^2}
\approx \qty{3.3e-7}{N}.
\]
Het \omterm{def:g10:universal-gravitation:weight}{gewicht} van elke vriend is ongeveer \qty{6.9e2}{N} --- twee
\emph{miljard} keer groter. Hun onderlinge aantrekking is echt, maar
hopeloos om te voelen.
\end{example}

\begin{example}[De Aarde houdt de Maan vast]\label{ex:g10:universal-gravitation:earthmoon}
Met $m_1 = \qty{5.97e24}{kg}$ (Aarde), $m_2 = \qty{7.35e22}{kg}$ (Maan)
en $d = \qty{3.84e8}{m}$:
\[
F = \num{6.67e-11} \times
\frac{\num{5.97e24} \times \num{7.35e22}}{(\num{3.84e8})^2}
\approx \qty{2.0e20}{N}.
\]
En de Maan trekt met dezelfde \qty{2.0e20}{N} aan de Aarde terug --- een
trek waarop de oceanen twee keer per dag antwoorden, als de getijden.
\end{example}

\begin{method}[Een gravitatiekracht berekenen]\label{met:g10:universal-gravitation:compute}
\begin{enumerate}
  \item Zet de massa's om naar \omterm{def:g10:orders-of-magnitude:unit}{kilogram} en de afstand tussen de
        middelpunten naar \omterm{def:g10:orders-of-magnitude:unit}{meter}.
  \item Pas $F = G m_1 m_2 / d^2$ toe.
  \item Controleer de macht van tien los van de cijfers, zoals in
        \cref{ch:g10:orders-of-magnitude}: \omterm{def:g10:universal-gravitation:force}{gravitatiekrachten} lopen van
        $\num{e-7}$ newton (twee mensen op een \omterm{def:g10:orders-of-magnitude:unit}{meter} afstand) tot
        $\num{e22}$ newton (Zon op Aarde), dus een verschoven \omterm{def:g10:orders-of-magnitude:scinot}{exponent}
        valt makkelijk op.
\end{enumerate}
\end{method}

\section{Gewicht en de sterkte van de zwaartekracht}

\begin{definition}[Gewicht]\label{def:g10:universal-gravitation:weight}
Het \emph{gewicht}\index{gewicht} $\vect P$ van een voorwerp aan het
oppervlak van een hemellichaam (een planeet, een maan, \dots) is de
\omterm{def:g10:universal-gravitation:force}{gravitatiekracht} die dat hemellichaam erop uitoefent. Het wijst naar het
middelpunt van het hemellichaam --- de richting die we
\emph{verticaal}\index{verticaal} noemen --- en wordt in newton gemeten.
\end{definition}

\begin{proposition}[Gewicht aan het oppervlak van een hemellichaam]\label{prop:g10:universal-gravitation:surfaceg}
Aan het oppervlak van een bolvormig hemellichaam met massa $M$ en straal
$R$ heeft een voorwerp met massa $m$ het \omterm{def:g10:universal-gravitation:weight}{gewicht}
\[
P = m\,g,
\qquad\text{waarin}\qquad
g = \frac{G M}{R^2}
\]
alleen van het hemellichaam afhangt, niet van het voorwerp.
\end{proposition}

\begin{proof}
Het middelpunt van het voorwerp en dat van het hemellichaam liggen een
afstand $R$ uit elkaar, dus geeft de wet van de algemene \omterm{def:g10:universal-gravitation:force}{gravitatie}
(\cref{thm:g10:universal-gravitation:law})
$P = G M m / R^2 = m \times \left(G M / R^2\right)$. De uitdrukking
tussen haakjes is voor elk voorwerp aan het oppervlak dezelfde: noem haar
$g$.
\end{proof}

\begin{definition}[Gravitatieveldsterkte]\label{def:g10:universal-gravitation:fieldstrength}
De grootheid $g = G M / R^2$, in newton per \omterm{def:g10:orders-of-magnitude:unit}{kilogram}, is de
\emph{gravitatieveldsterkte}\index{gravitatieveldsterkte} aan het
oppervlak van het hemellichaam: het trekt aan elke \omterm{def:g10:orders-of-magnitude:unit}{kilogram} met $g$
newton. Op Aarde is $g = \qty{9.81}{N/kg}$.
\end{definition}

\begin{example}[$g$ berekenen]\label{ex:g10:universal-gravitation:gvalues}
Voor de Aarde ($M = \qty{5.97e24}{kg}$, $R = \qty{6.37e6}{m}$):
\[
g = \frac{\num{6.67e-11} \times \num{5.97e24}}{(\num{6.37e6})^2}
= \frac{\num{3.98e14}}{\num{4.06e13}}
\approx \qty{9.81}{N/kg}.
\]
Voor de Maan ($M = \qty{7.35e22}{kg}$, $R = \qty{1.74e6}{m}$) geeft
dezelfde berekening $g_{\text{Maan}} \approx \qty{1.62}{N/kg}$ --- zes
keer minder: de Maan is veel lichter, en kleiner zijn maakt dat niet
goed. Voor Mars is $g_{\text{Mars}} \approx \qty{3.73}{N/kg}$.
\end{example}

\begin{omfigure}
\begin{tikzpicture}
\begin{axis}[omaxis, axis x line=bottom, axis y line=left,
  width=8.8cm, height=5.4cm, ybar, bar width=15pt,
  symbolic x coords={Maan, Mars, Aarde, Jupiter}, xtick=data,
  ylabel={$g$ (\unit{N/kg})}, ymin=0, ymax=29,
  nodes near coords,
  every node near coord/.append style={font=\small},
  enlarge x limits=0.16]
\addplot[draw=omDef, fill=omDef!25] coordinates
  {(Maan,1.62) (Mars,3.73) (Aarde,9.81) (Jupiter,24.8)};
\end{axis}
\end{tikzpicture}

{\small \omterm{def:g10:universal-gravitation:fieldstrength}{Gravitatieveldsterkte} aan het oppervlak van vier werelden:
dezelfde zak suiker van \qty{1}{kg} weegt \qty{1.6}{N} op de Maan en
\qty{25}{N} op Jupiter.}
\end{omfigure}

\begin{omfigure}
\begin{tikzpicture}[scale=1.0]
  \fill[omExo!15] (0,0) circle (1.5);
  \draw[omExo, thick] (0,0) circle (1.5);
  \fill (0,0) circle (1.5pt) node[below, font=\small] {$O$};
  \fill (90:1.5) circle (1.8pt) node[above, font=\small] {$A$};
  \draw[-Stealth, omThm, thick] (90:1.5) -- (90:0.8)
    node[pos=0.6, right, font=\small] {$\vect P_A$};
  \fill (200:1.5) circle (1.8pt) node[left, font=\small] {$B$};
  \draw[-Stealth, omThm, thick] (200:1.5) -- (200:0.8)
    node[pos=0.6, below, font=\small] {$\vect P_B$};
  \fill (330:1.5) circle (1.8pt) node[right, font=\small] {$C$};
  \draw[-Stealth, omThm, thick] (330:1.5) -- (330:0.8)
    node[pos=0.6, above left, font=\small] {$\vect P_C$};
\end{tikzpicture}

{\small ``Omlaag'' is een plaatselijk begrip: in elk punt van de
aardbol wijst de gewichtsvector naar het middelpunt $O$. Twee wandelaars
op tegenovergestelde plekken staan allebei rechtop --- met de voeten naar
elkaar toe.}
\end{omfigure}

\begin{remark}[Waarom astronauten stuiteren]\label{rem:g10:universal-gravitation:bounce}
Een astronaut met een massa van \qty{120}{kg} (lichaam plus pak) weegt op
Aarde $120 \times 9.81 \approx \qty{1180}{N}$, maar op de Maan slechts
$120 \times 1.62 \approx \qty{194}{N}$. Spieren die onder \qty{1180}{N}
zijn getraind, dragen plots een zesde van de last: vandaar de beroemde
trage, huppelende maangang. De massa --- en daarmee de moeite die het
kost om te \emph{stoppen} of te \emph{draaien} --- verandert niet, en
daarom vielen de maanwandelaars zo vaak om.
\end{remark}

\section{Massa of gewicht? De balans en de veerunster}

In het dagelijks taalgebruik ``weegt'' een zak meel drie kilo. De
natuurkunde splitst die zin in tweeën.

\begin{remark}[Massa is geen gewicht]\label{rem:g10:universal-gravitation:massweight}
\begin{itemize}
  \item De \emph{massa} $m$ meet de hoeveelheid materie, in \omterm{def:g10:orders-of-magnitude:unit}{kilogram}. Ze
        is een eigenschap van het voorwerp alleen: dezelfde op Aarde, op
        de Maan of ergens in de diepe ruimte.
  \item Het \emph{\omterm{def:g10:universal-gravitation:weight}{gewicht}} $P = mg$ is een kracht, in newton, die
        \emph{door een hemellichaam} op het voorwerp wordt uitgeoefend.
        Het verandert van hemellichaam tot hemellichaam met $g$ --- en
        het verdwijnt ver van elk hemellichaam.
\end{itemize}
De twee verwarren is op de markt onschuldig en in een natuurkundige
opgave fataal: de \omterm{def:g10:orders-of-magnitude:unit}{kilogram} is geen \omterm{def:g10:orders-of-magnitude:unit}{eenheid} van kracht.
\end{remark}

\begin{method}[De balans en de veerunster]\label{met:g10:universal-gravitation:instruments}
\begin{itemize}
  \item Een \emph{balans} vergelijkt het voorwerp met standaardmassa's:
        de \omterm{def:g10:universal-gravitation:weight}{gewichten} in beide schalen worden met dezelfde $g$
        geschaald, dus meet de vergelijking de \emph{massa} --- op elk
        hemellichaam hetzelfde oordeel.
  \item Een \emph{veerunster} meet de kracht die haar veer uitrekt: ze
        leest het \emph{\omterm{def:g10:universal-gravitation:weight}{gewicht}} af, in newton, en haar aanwijzing
        verandert van hemellichaam tot hemellichaam. Om de massa terug
        te vinden, deel je door de plaatselijke veldsterkte:
        $m = P/g$.
\end{itemize}
\end{method}

\begin{example}[Meel op de Maan]\label{ex:g10:universal-gravitation:flour}
Een zak meel van \qty{3.0}{kg} houdt \qty{3.0}{kg} standaardmassa's in
evenwicht op Aarde \emph{en} op de Maan. Een veerunster wijst echter
$3.0 \times 9.81 \approx \qty{29}{N}$ aan op Aarde en
$3.0 \times 1.62 \approx \qty{4.9}{N}$ op de Maan. Een handelaar die op
de Maan ``per newton'' verkoopt, zou zes keer zoveel meel meegeven ---
de handel draait, net als de wetenschap, op massa.
\end{example}

\section{Om de Aarde heen vallen: het kanon van Newton}

Gooi een steen horizontaal weg: de zwaartekracht buigt zijn baan tot een
boog en hij landt een paar \omterm{def:g10:orders-of-magnitude:unit}{meter} verderop. Gooi harder: hij landt verder.
Newtons gedachte-experiment drijft dat idee tot het uiterste. De Aarde is
rond, en haar oppervlak zakt ongeveer \qty{5}{m} onder de horizontaal per
\qty{8}{km} afgelegde weg; een vrij vallend lichaam valt in zijn eerste
seconde \qty{4.9}{m} (\cref{ch:g12:projectile-motion} zal dat vaststellen).
Een projectiel dat met \qty{8}{km/s} langs de Aarde scheert, valt dus
precies even snel als de grond wegkromt: het valt eeuwig zonder ooit te
landen.

\begin{definition}[Satelliet en omloopbaan]\label{def:g10:universal-gravitation:satellite}
Een \emph{satelliet}\index{satelliet} van een hemellichaam is een lichaam
dat om dat hemellichaam heen valt zonder het oppervlak ooit te bereiken;
de weg die het steeds herhaalt, is zijn \emph{omloopbaan}\index{omloopbaan}.
\end{definition}

\begin{omfigure}
\begin{tikzpicture}[scale=1.0]
  \fill[omExo!15] (0,0) circle (1.5);
  \draw[omExo, thick] (0,0) circle (1.5);
  \fill[omExo!60] (-0.22,1.44) -- (0,1.85) -- (0.22,1.44) -- cycle;
  \coordinate (L) at (0,1.85);
  \draw[omThm, thick] (L) .. controls (0.6,1.85) .. (55:1.5);
  \draw[omProp, thick] (L) .. controls (1.1,1.85) and (1.7,1.1) .. (10:1.5);
  \draw[omMeth, thick] (L) .. controls (1.8,1.8) and (2.3,-0.4) .. (-70:1.5);
  \draw[-Stealth, omDef, thick] (0,1.85) arc (90:-250:1.85);
  \node[font=\small, omDef, anchor=east] at (-1.5,1.45) {omloopbaan};
\end{tikzpicture}

{\small Het kanon van Newton: steeds sneller afgevuurd vanaf een
bergtop, landt de kogel steeds verder --- tot bij ongeveer
\qty{7.9}{km/s} de grond even snel wegkromt als de kogel valt, en de
kogel in een \omterm{def:g10:universal-gravitation:satellite}{omloopbaan} komt.}
\end{omfigure}

\begin{remark}[De Maan is een vallende steen]\label{rem:g10:universal-gravitation:fallingmoon}
De Maan is de natuurlijke \omterm{def:g10:universal-gravitation:satellite}{satelliet} van de Aarde: ze valt elke seconde
ongeveer \qty{1.4}{mm} naar ons toe, en haar zijwaartse beweging voert
haar om ons heen voordat ze dichterbij kan komen
(in \cref{exo:g10:universal-gravitation:15} doe je Newtons eigen controle
van dat getal over). De astronauten van een ruimtestation in een \omterm{def:g10:universal-gravitation:satellite}{omloopbaan}
zitten in dezelfde vrije val als hun schip --- daarom zweven ze, en niet
doordat de zwaartekracht zou ontbreken. De \omterm{def:g10:universal-gravitation:force}{gravitatie} vormt elke baan,
van de boog van de steen tot de ellips van de planeet; het kwantitatieve
verhaal van de omloopbanen staat in \cref{ch:g12:satellites-kepler}.
\end{remark}

\begin{notation}[Gegevenskaart]
Neem, tenzij een oefening iets anders zegt,
$G = \qty{6.67e-11}{N.m^2/kg^2}$ en:
\begin{center}
\small
\begin{tabular}{lccc}
 & massa (\unit{kg}) & straal (\unit{m}) & $g$ (\unit{N/kg}) \\
Aarde & \num{5.97e24} & \num{6.37e6} & 9.81 \\
Maan & \num{7.35e22} & \num{1.74e6} & 1.62 \\
Mars & \num{6.42e23} & \num{3.39e6} & 3.73 \\
Zon & \num{1.99e30} & --- & --- \\
\end{tabular}
\end{center}
Afstand Aarde--Maan: \qty{3.84e8}{m}; afstand Aarde--Zon:
\qty{1.496e11}{m} (van middelpunt tot middelpunt).
\end{notation}

\section{Oefeningen}

\begin{exercise}[$\star$]\label{exo:g10:universal-gravitation:1}
Twee danspartners van \qty{60}{kg} houden elkaar vast met hun
middelpunten \qty{0.80}{m} uit elkaar.
\begin{enumerate}
  \item Bereken hun onderlinge gravitatieaantrekking.
  \item Vergelijk die met het \omterm{def:g10:universal-gravitation:weight}{gewicht} van een mug van \qty{2.5}{mg}.
\end{enumerate}
\end{exercise}

\begin{exercise}[$\star$]\label{exo:g10:universal-gravitation:2}
Een leerlinge heeft een massa van \qty{55}{kg}.
\begin{enumerate}
  \item Bereken haar \omterm{def:g10:universal-gravitation:weight}{gewicht} op Aarde.
  \item Ze reist naar de Maan. Wat zijn daar haar massa en haar \omterm{def:g10:universal-gravitation:weight}{gewicht}?
\end{enumerate}
\end{exercise}

\begin{exercise}[$\star$]\label{exo:g10:universal-gravitation:3}
\begin{enumerate}
  \item Bereken uit de gegevenskaart de \omterm{def:g10:universal-gravitation:fieldstrength}{gravitatieveldsterkte} aan het
        oppervlak van Mars opnieuw.
  \item Er wordt een wagentje van \qty{900}{kg} heen gestuurd. Bereken
        zijn \omterm{def:g10:universal-gravitation:weight}{gewicht} op Mars en zijn \omterm{def:g10:universal-gravitation:weight}{gewicht} op Aarde.
\end{enumerate}
\end{exercise}

\begin{exercise}[$\star$]\label{exo:g10:universal-gravitation:4}
Bereken met de gegevenskaart de kracht die de Aarde op de Maan
uitoefent. Welke kracht oefent de Maan op de Aarde uit?
\end{exercise}

\begin{exercise}[$\star$]\label{exo:g10:universal-gravitation:5}
Twee lichamen trekken elkaar aan met een kracht $F_0$. Wat wordt de
kracht als:
\begin{enumerate}
  \item de afstand tussen de middelpunten verdubbelt;
  \item de afstand halveert;
  \item beide massa's verdubbelen (bij dezelfde afstand);
  \item één massa verdrievoudigt \emph{én} de afstand verdrievoudigt?
\end{enumerate}
\end{exercise}

\begin{exercise}[$\star\star$]\label{exo:g10:universal-gravitation:6}
Twee geladen supertankers van elk \qty{5.0e8}{kg} liggen afgemeerd met
hun middelpunten \qty{100}{m} uit elkaar.
\begin{enumerate}
  \item Bereken hun onderlinge gravitatieaantrekking.
  \item Welke massa heeft op Aarde een \omterm{def:g10:universal-gravitation:weight}{gewicht} dat gelijk is aan die
        kracht?
  \item De kracht is verrassend groot --- en toch vertonen de tankers
        geen neiging naar elkaar toe te drijven. Bereken de versnelling
        $a = F/m$ die ze aan één tanker geeft, en becommentarieer.
\end{enumerate}
\end{exercise}

\begin{exercise}[$\star\star$]\label{exo:g10:universal-gravitation:7}
Een winkel op de Maan verkoopt een zak rijst van \qty{3.0}{kg}.
\begin{enumerate}
  \item Wat wijzen een balans en een veerunster voor deze zak aan op
        Aarde? En op de Maan?
  \item Een klant betaalt voor ``\qty{3.0}{kg}''. Welk instrument
        waarborgt een eerlijke koop op elke wereld, en waarom?
\end{enumerate}
\end{exercise}

\begin{exercise}[$\star\star$]\label{exo:g10:universal-gravitation:8}
Het internationale ruimtestation vliegt op een hoogte van \qty{400}{km}.
\begin{enumerate}
  \item Bereken de \omterm{def:g10:universal-gravitation:fieldstrength}{gravitatieveldsterkte} van de Aarde op die hoogte.
        (Let op: wat is de afstand tot het \emph{middelpunt} van de
        Aarde?)
  \item Druk die uit als percentage van de waarde aan het oppervlak.
  \item De astronauten aan boord zweven. Verzoen dat met je antwoord.
\end{enumerate}
\end{exercise}

\begin{exercise}[$\star\star$]\label{exo:g10:universal-gravitation:9}
Aan het oppervlak van Venus is $g = \qty{8.87}{N/kg}$, en de straal van
de planeet is \qty{6.05e6}{m}. Leid de massa van Venus af en vergelijk
die met die van de Aarde.
\end{exercise}

\begin{exercise}[$\star\star$]\label{exo:g10:universal-gravitation:10}
\begin{enumerate}
  \item Bereken de kracht die de Zon op \qty{1.0}{kg} zeewater
        uitoefent, en daarna de kracht die de Maan op dezelfde \omterm{def:g10:orders-of-magnitude:unit}{kilogram}
        uitoefent.
  \item De Zon trekt harder --- met welke factor?
  \item En toch beheerst de Maan de getijden. Het getij wordt niet
        gedreven door de trek zelf maar door het \emph{verschil} tussen
        de trekkracht aan de dichtstbijzijnde en aan de verste kant van
        de Aarde. Leg zonder te rekenen uit waarom de nabije Maan het op
        dat punt van de verre Zon kan winnen.
\end{enumerate}
\end{exercise}

\begin{exercise}[$\star\star$]\label{exo:g10:universal-gravitation:11}
\begin{enumerate}
  \item Een planeet heeft tweemaal de massa van de Aarde en tweemaal
        haar straal. Druk haar veldsterkte aan het oppervlak uit in de
        $g$ van de Aarde, en daarna in \unit{N/kg}.
  \item Welke straal zou een planeet met de massa van de Aarde nodig
        hebben opdat haar veldsterkte aan het oppervlak $2g$ is?
\end{enumerate}
\end{exercise}

\begin{exercise}[$\star\star\star$]\label{exo:g10:universal-gravitation:12}
Ergens op de lijn Aarde--Maan heffen de twee trekkrachten op een
ruimtesonde elkaar op. Zij $x$ de afstand van het middelpunt van de
Aarde tot dat punt en $D$ de afstand Aarde--Maan.
\begin{enumerate}
  \item Toon aan dat in dat punt
        $\left(\dfrac{x}{D - x}\right)^{2}
        = \dfrac{M_{\text{Aarde}}}{M_{\text{Maan}}}$.
  \item Leid $x$ af. Welk deel van de reis naar de Maan is dat?
\end{enumerate}
\end{exercise}

\begin{exercise}[$\star\star\star$]\label{exo:g10:universal-gravitation:13}
Het kanon van Newton, met getallen. Een projectiel scheert horizontaal
met snelheid $v$ langs de Aarde.
\begin{enumerate}
  \item Na $x = \qty{8.0}{km}$ horizontale weg is de bolvormige Aarde
        ``weggezakt'' over een hoogte $h$ die voldoet aan
        $(R + h)^2 = R^2 + x^2$. Toon met
        $(R + h)^2 \approx R^2 + 2Rh$ voor kleine $h$ aan dat
        $h \approx x^2 / (2R)$, en bereken $h$.
  \item Een vrij vallend lichaam valt in zijn eerste seconde
        \qty{4.9}{m}. Vergelijk met $h$ en leid de snelheid af waarbij
        het projectiel nooit landt.
  \item Bereken de duur van één volledige omloop bij die snelheid, vlak
        langs het oppervlak.
\end{enumerate}
\end{exercise}

\begin{exercise}[$\star\star\star$]\label{exo:g10:universal-gravitation:14}
Een neutronenster propt $M = \qty{2.8e30}{kg}$ (1.4 maal de massa van de
Zon) in een straal van \qty{12}{km}.
\begin{enumerate}
  \item Bereken de \omterm{def:g10:universal-gravitation:fieldstrength}{gravitatieveldsterkte} aan haar oppervlak.
  \item Wat zou een mens van \qty{70}{kg} daar wegen? Vergelijk met zijn
        \omterm{def:g10:universal-gravitation:weight}{gewicht} op Aarde.
  \item Wat is daar het \omterm{def:g10:universal-gravitation:weight}{gewicht} van een zandkorrel van \qty{1.0}{mg}?
        Welke massa heeft dat \omterm{def:g10:universal-gravitation:weight}{gewicht} op Aarde?
\end{enumerate}
\end{exercise}

\begin{exercise}[$\star\star\star$]\label{exo:g10:universal-gravitation:15}
Newtons maantoets (1666). De straal van de maanbaan is ongeveer $60$
aardstralen.
\begin{enumerate}
  \item Met welke factor is de trek van de Aarde per \omterm{def:g10:orders-of-magnitude:unit}{kilogram} zwakker bij
        de Maan dan aan het aardoppervlak, als de zwaartekracht het
        omgekeerd kwadratische verband volgt?
  \item Vlak bij het oppervlak valt een lichaam in zijn eerste seconde
        \qty{4.9}{m}. Hoever zou de Maan elke seconde naar de Aarde toe
        moeten vallen?
  \item Toets dat aan de \omterm{def:g10:universal-gravitation:satellite}{omloopbaan}: de Maan doorloopt haar cirkel met straal
        $d = \qty{3.84e8}{m}$ in $T = 27.3$ dagen. Bereken haar snelheid
        $v$, de afstand $x$ die ze in één seconde aflegt, en daarna de
        val $h \approx x^2/(2d)$ (zoals in
        \cref{exo:g10:universal-gravitation:13}).
  \item Vergelijk de antwoorden van vraag 2 en 3, en zeg wat Newton
        eruit concludeerde.
\end{enumerate}
\end{exercise}

\section{Probleem: De Aarde wegen}

\begin{problem}[{Weekendopgave --- van twee loden bollen in een schuurtje
naar de massa van de Aarde en daarna van de Zon: wat één kleine
constante, eenmaal gemeten, waard is}]
\label{pb:g10:universal-gravitation:1}
In 1798 hing Henry Cavendish een lichte staaf met twee kleine loden
bollen aan een dunne draad, bracht twee grote loden bollen dichtbij en
mat de bijna onbestaande verdraaiing van de draad. De kranten schreven
dat hij \emph{de Aarde had gewogen} --- en ze hadden gelijk: zodra $G$
bekend is, geven de $g$ en de straal van de Aarde haar massa prijs, en
geeft de jaarlijkse \omterm{def:g10:universal-gravitation:satellite}{omloopbaan} van de Aarde die van de Zon prijs. Deze opgave
loopt de hele keten na. Gebruik overal de gegevenskaart.

\medskip
\noindent\textbf{Deel I --- De zwakste kracht.}
\begin{enumerate}
  \item Het ijzer van de Eiffeltoren heeft een massa van ongeveer
        \qty{7.3e6}{kg}. Een toerist van \qty{55}{kg} staat \qty{100}{m}
        van het massamiddelpunt ervan. Bereken de trek van de toren op de
        toerist.
  \item Vergelijk die met het \omterm{def:g10:universal-gravitation:weight}{gewicht} van een zandkorrel van
        \qty{1.0}{mg}.
  \item Oefening in evenredigheid: wat gebeurt er met een
        \omterm{def:g10:universal-gravitation:force}{gravitatiekracht} als (a) de afstand met $10$ wordt
        vermenigvuldigd; (b) beide massa's met $1000$ worden
        vermenigvuldigd; (c) beide massa's \emph{en} de afstand met
        $1000$ worden vermenigvuldigd?
  \item Bereken uit $M$, $R$ en $G$ de trek van de Aarde op een voorwerp
        van \qty{1.00}{kg} aan haar oppervlak. Welke alledaagse naam en
        welk symbool draagt dit getal?
  \item De zwaartekracht is veruit de zwakste wisselwerking die je bent
        tegengekomen --- en toch regeert ze het heelal. Leg in een of
        twee zinnen uit waarom de trek van de Aarde op jou die van je
        buurman overheerst.
\end{enumerate}

\medskip
\noindent\textbf{Deel II --- De balans van Cavendish.}
In de torsiebalans trekt een grote bol met massa $m_1 = \qty{158}{kg}$
aan een kleine met massa $m_2 = \qty{0.73}{kg}$; hun middelpunten liggen
$d = \qty{0.23}{m}$ uit elkaar.
\begin{enumerate}[resume]
  \item Bereken met de moderne waarde van $G$ de kracht tussen de twee
        bollen.
  \item Bereken het \omterm{def:g10:universal-gravitation:weight}{gewicht} van de kleine bol en de verhouding van de
        twee krachten. Waarom had Cavendish een torsiedraad nodig en niet
        een gewone weegschaal?
  \item De verdraaiing die Cavendish mat, komt (in moderne eenheden)
        overeen met $F = \qty{1.47e-7}{N}$. Leid zijn waarde van $G$ af
        en vergelijk die met de huidige.
  \item Leid met $G = \qty{6.67e-11}{N.m^2/kg^2}$, $g = \qty{9.81}{N/kg}$
        en $R = \qty{6.37e6}{m}$ de massa van de Aarde af. (Dit is de
        stap die de kranten vierden.)
  \item Leid de gemiddelde dichtheid $\rho = M / V$ van de Aarde af, met
        $V = \frac43 \pi R^3$. Gesteente aan het oppervlak heeft
        dichtheden rond \qty{2.7e3}{kg/m^3}: wat onthult de vergelijking
        over het diepe binnenste?
\end{enumerate}

\medskip
\noindent\textbf{Deel III --- Andere werelden.}
\begin{enumerate}[resume]
  \item Bereken uit de gegevenskaart de veldsterkte aan het oppervlak van
        de Maan opnieuw.
  \item Op Aarde tilt een goede afzet je massamiddelpunt \qty{0.50}{m}
        omhoog. Bij dezelfde afzet is de bereikte hoogte omgekeerd
        evenredig met $g$ (zoals het energiehoofdstuk,
        \cref{ch:g10:energy-conservation}, zal verantwoorden). Hoe hoog
        brengt dezelfde sprong je op de Maan?
  \item Mercurius: $M = \qty{3.30e23}{kg}$, $R = \qty{2.44e6}{m}$.
        Bereken zijn veldsterkte aan het oppervlak en vergelijk die met
        die van Mars. Mercurius is veel kleiner dan Mars --- hoe kunnen
        de twee waarden dan zo dicht bij elkaar liggen?
  \item Op welke hoogte boven het aardoppervlak is de veldsterkte tot de
        helft van de waarde aan het oppervlak gezakt?
  \item Op welke afstand van het middelpunt van de Aarde daalt de trek
        van de Aarde per \omterm{def:g10:orders-of-magnitude:unit}{kilogram} tot de waarde aan het maanoppervlak,
        \qty{1.62}{N/kg}? Druk die uit in aardstralen.
\end{enumerate}

\medskip
\noindent\textbf{Deel IV --- De Zon wegen.}
De Aarde doorloopt in één jaar, $T = \qty{3.156e7}{s}$, een vrijwel
cirkelvormige \omterm{def:g10:universal-gravitation:satellite}{omloopbaan} met straal $d = \qty{1.496e11}{m}$. Neem als gegeven
(het wordt bewezen in \cref{ch:g12:satellites-kepler}) dat een lichaam op
een cirkelvormige baan met straal $d$ en snelheid $v$ een trek van
$v^2/d$ newton per \omterm{def:g10:orders-of-magnitude:unit}{kilogram} naar het middelpunt nodig heeft.
\begin{enumerate}[resume]
  \item Bereken de baansnelheid $v$ van de Aarde.
  \item Leid daaruit de trek per \omterm{def:g10:orders-of-magnitude:unit}{kilogram} af die de Zon op de Aarde
        uitoefent.
  \item Die trek is ook $G M_{\text{Zon}} / d^2$. Leid de massa van de
        Zon af.
  \item Hoeveel Aardes is dat? Vergelijk $M_{\text{Zon}}$ met
        $M_{\text{Aarde}}$.
  \item Slot. De bollen van Cavendish, de Aarde, de Zon: één constante
        $G$ volstond voor alle drie. Zeg in één zin wat het woord
        \emph{algemene} in ``algemene \omterm{def:g10:universal-gravitation:force}{gravitatie}'' ons in deze opgave
        heeft opgeleverd.
\end{enumerate}
\end{problem}
